\part{Reference}

\chapter{Comfort Presets and Visibility}
\label{ch:comfort-presets}

\rowcolors{2}{LtGray}{white}
\begin{longtable}{C{1.4cm}L{2.0cm}L{4.0cm}L{6.7cm}}
\toprule
\rowcolor{Navy}
\color{white}\sffamily\bfseries Icon &
\color{white}\sffamily\bfseries Preset &
\color{white}\sffamily\bfseries Conditions &
\color{white}\sffamily\bfseries Key characteristics \\
\midrule
\iconlg{comfort-default} & Default & General fallback &
  Balanced mid-contrast; usable in most conditions. \\
\iconlg{comfort-dawn}    & Dawn    & Around sunrise &
  Warm tones, slightly reduced brightness. \\
\iconlg{comfort-day}     & Day     & Full daylight, direct sun &
  Maximum brightness and contrast. For open-ocean sailing. \\
\iconlg{comfort-sunset}  & Sunset  & Around sunset &
  Warmer, transitional brightness. \\
\iconlg{comfort-dusk}    & Dusk    & Twilight &
  Reduced brightness; begins night-vision preservation. \\
\iconlg{comfort-night}   & Night   & Full dark, offshore watches &
  Minimum brightness, dark backgrounds, red-tinted icons.
  Preserves night vision.
  Alarms remain high-contrast on every preset. \\
\bottomrule
\end{longtable}

\section{Practical Guidance}

\begin{itemize}
  \item Use brighter, higher-contrast presets in daylight.
  \item Use dimmer, night-oriented presets when preserving night vision matters.
  \item Keep alarms and safety-critical text easy to read in every preset.
  \item Prefer readability over appearance. If a preset looks attractive but makes
        urgent information harder to see, switch to a clearer one.
  \item The Night preset sets the icon tint to red by default.
        This is intentional --- red light minimises the time needed to
        re-adapt to darkness after looking at the display.
\end{itemize}

\chapter{Display, Touch, and Rough-Conditions Operation}
\label{ch:roughconditions}

\section{Display Readability}

FairWindSK should remain readable in bright daylight, dusk, night, and rough
conditions.
If the display is difficult to read, correct that first before diagnosing any
app issues.

Start with:
\begin{enumerate}
  \item Screen brightness.
  \item Comfort preset (correct for the current light conditions).
  \item Viewing angle and helm position.
  \item The app or page shown in the Application Area.
\end{enumerate}

Keep safety-critical numbers easy to identify at a glance.
Avoid presets that reduce contrast too much.
Use larger displays for primary navigation whenever possible.
Use a display sunshade or hood in direct sunlight.

\section{Touch Use in Rough Conditions}

\begin{itemize}
  \item Brace yourself first. One hand on the boat, one on the screen.
  \item Use deliberate, single taps. Hesitant pressure and repeated rapid taps
        cause unintended actions.
  \item Never explore the autopilot controls by tapping speculatively ---
        every command has an immediate physical effect.
  \item Keep the screen clean. Salt spray reduces touch accuracy significantly.
  \item Pre-configure what you need while conditions are calm.
  \item Prefer the large, stable targets of the Bottom Bar main icons.
  \item Confirm dialogs and mode changes before acting on them.
\end{itemize}

For helm use, stable layout and good readability matter more than decorative styling.

\chapter{Recommended Operator Workflow}
\label{ch:workflow}

For most passages, this sequence keeps the shell predictable and easy to trust.

\section{Before Departure}

\begin{enumerate}
  \item Start FairWindSK and confirm the shell health state is normal.
  \item Verify that expected key data is live in the Top Bar.
  \item Open the app most relevant to the current passage phase.
  \item Confirm brightness and comfort preset.
  \item Complete the pre-departure checklist (Chapter~2).
\end{enumerate}

\section{During Passage}

\begin{enumerate}
  \item Watch for stale or missing values in the Top Bar.
  \item Cross-check route, depth, and heading whenever conditions change.
  \item Use the launcher to switch tasks rather than forcing a stuck app.
  \item Treat degraded states as operational warnings.
  \item If the current screen stops feeling clear or trustworthy, tap
        \iconlg{applications}~\key{Apps} to return to the launcher.
\end{enumerate}

\section{After a Reconnect or Server Restart}

\begin{enumerate}
  \item Wait briefly for automatic recovery to complete.
  \item Confirm the status shows \textit{Live}.
  \item Reopen the app if the foreground application was degraded.
  \item Reconfirm key navigation values before depending on them.
\end{enumerate}

\section{Task Switching}

Use the Bottom Bar rather than forcing the current app to do everything:

\begin{itemize}
  \item Use \iconlg{applications}~Apps to return to the launcher and switch tasks.
  \item Use the open-apps strip to jump directly to a loaded application.
  \item Use operational bars (\iconlg{alarm-pob}~POB, \iconlg{anchor-iec}~Anchor,
        \iconlg{command-autopilot}~Autopilot) for vessel-operation tasks
        without leaving your navigation app.
\end{itemize}

\chapter{Signal K Server on Docker}
\label{ch:docker}

\section{Official Docker Image}

\begin{lstlisting}
cr.signalk.io/signalk/signalk-server:latest
\end{lstlisting}

Supported architectures: \code{linux/amd64}, \code{linux/arm/v7},
\code{linux/arm64} (all Raspberry~Pi~4 and~5 configurations).

\section{Quick Start with Demo NMEA Data}

\begin{lstlisting}
mkdir $HOME/.signalk
docker run --init -it --rm \
  --name signalk-server \
  --publish 3000:3000 \
  -v $HOME/.signalk:/home/node/.signalk \
  cr.signalk.io/signalk/signalk-server \
  --sample-nmea0183-data
\end{lstlisting}

\section{Production Deployment}

\begin{lstlisting}
docker run -d --init \
  --name signalk-server \
  -p 3000:3000 \
  -v $HOME/.signalk:/home/node/.signalk \
  cr.signalk.io/signalk/signalk-server
\end{lstlisting}

\section{Updating Signal K}

\begin{lstlisting}
docker pull cr.signalk.io/signalk/signalk-server:latest
docker stop signalk-server
docker rm signalk-server
# Re-run the production docker run command above
\end{lstlisting}

\chapter{Troubleshooting}
\label{ch:troubleshooting}

\section{I See Missing Values}

Possible causes: the sensor or source instrument is offline; Signal~K is not
receiving that path; the value is not configured for display; the connection has
dropped.

\begin{enumerate}
  \item Check shell health state.
  \item See whether other values are still live.
  \item Confirm the source instrument is operating.
  \item Check \ui{Settings~> Signal~K Paths} --- is the correct path configured?
  \item Return to the launcher and reopen the affected app if needed.
\end{enumerate}

\section{I See Stale Values}

Possible causes: temporary data interruption; reconnect in progress; network delay;
partial degradation in the data source.

\begin{enumerate}
  \item Treat the values cautiously.
  \item Wait briefly for recovery.
  \item Cross-check with other instruments.
  \item Investigate if the stale state persists beyond 60 seconds.
\end{enumerate}

\section{A Web App Is Blank or Frozen}

\begin{enumerate}
  \item Retry once if the shell offers the option.
  \item Return to the launcher.
  \item Check shell health and Signal~K connection state.
  \item Use another app if the shell is otherwise healthy.
  \item A failed application does not mean FairWindSK has failed.
\end{enumerate}

\section{The Shell Is Disconnected}

\begin{enumerate}
  \item Is Signal~K running? Check LEDs on the Raspberry~Pi; check the Docker
        container status: \code{docker ps -a | grep signalk}.
  \item Is your device on the correct Wi-Fi network?
  \item Is the server address in \ui{Settings~> Connection} correct?
  \item Can a browser on another device reach \code{http://server-address:3000}?
  \item Allow FairWindSK 60 seconds to attempt automatic reconnection.
  \item Reconfirm all safety-critical values after reconnection.
\end{enumerate}

\section{No Applications Appear on the Launcher}

\begin{enumerate}
  \item Confirm status shows \textit{Live}.
  \item Are apps installed in the Signal~K App Store?
  \item In \ui{Settings~> Connection}, tap \key{Connect} to force a reconnect and
        app discovery.
\end{enumerate}

\section{Autopilot or Anchor Icon Is Greyed Out}

\begin{enumerate}
  \item Is the required plugin installed in the Signal~K App Store?
  \item Is the plugin enabled and configured in Signal~K settings?
  \item Is the correct plugin action path configured in
        \ui{Settings~> Signal~K Paths}?
\end{enumerate}

\section{FairWindSK Running Slowly on Raspberry Pi}

\begin{enumerate}
  \item Check CPU temperature: \code{vcgencmd measure\_temp}.
        Above 80°C indicates thermal throttling --- improve cooling.
  \item Check \ui{Settings~> System~> Performance} for memory usage.
        FairWindSK should use under 250\,MB at idle.
  \item Use a Class~A2 high-speed microSD card.
        A slow card is the most common cause of sluggish performance on Raspberry~Pi.
  \item Set Log Level to \textit{Off} in \ui{Settings~> System~> Diagnostics}
        to minimise SD card writes.
\end{enumerate}

\chapter{Daily Operating Checklist}
\label{ch:checklist}

\section*{Before Departure}

\rowcolors{2}{LtGray}{white}
\begin{longtable}{C{0.8cm}L{12.9cm}}
\toprule
\rowcolor{Navy}
\color{white}\sffamily\bfseries $\square$ &
\color{white}\sffamily\bfseries Check \\
\midrule
$\square$ & FairWindSK started --- status shows \textit{Live} \\
$\square$ & Position live and plausible \\
$\square$ & SOG updating --- valid GPS fix \\
$\square$ & Depth updating (if fitted) \\
$\square$ & Heading showing \\
$\square$ & Chart plotter open; vessel in correct position \\
$\square$ & Route or waypoint active (if proceeding to a destination) \\
$\square$ & AIS targets visible (if AIS fitted) \\
$\square$ & Comfort preset set for current light conditions \\
$\square$ & Screen brightness readable from helm \\
$\square$ & POB button location confirmed \\
$\square$ & Autopilot available and responding (if required) \\
$\square$ & Anchor watch available (if anchoring tonight) \\
\bottomrule
\end{longtable}

\section*{After Arrival}

\rowcolors{2}{LtGray}{white}
\begin{longtable}{C{0.8cm}L{12.9cm}}
\toprule
\rowcolor{Navy}
\color{white}\sffamily\bfseries $\square$ &
\color{white}\sffamily\bfseries Check \\
\midrule
$\square$ & Return to launcher --- no active alarms \\
$\square$ & Anchor watch set with appropriate radius (if anchoring) \\
$\square$ & Comfort preset changed to Dusk or Night \\
$\square$ & Track exported for logbook (if required) \\
\bottomrule
\end{longtable}

\chapter{Quick Reference Card}
\label{ch:quickref}

\section*{Immediate Action Guide}

\rowcolors{2}{LtGray}{white}
\begin{longtable}{L{3.6cm}L{4.4cm}L{5.7cm}}
\toprule
\rowcolor{Navy}
\color{white}\sffamily\bfseries Situation &
\color{white}\sffamily\bfseries Immediate action &
\color{white}\sffamily\bfseries Then \\
\midrule
Person Over Board &
  Tap \iconlg{alarm-pob}~\key{POB}. Once. &
  Call MAYDAY on VHF Ch~16. Deploy SAR route. \\
Alarm flashing &
  Tap \iconlg{alarm-general}~\key{Alarms}. Read it. &
  Act per emergency procedures. \\
Screen frozen &
  Wait 10~s. Tap \iconlg{applications}~\key{Apps}. &
  Check status. Try reopening the application. \\
Shell disconnected &
  Wait 60~s for auto-reconnect. &
  If still disconnected: check Signal~K and Wi-Fi. \\
Stale depth in shoal water &
  Cross-check immediately with a backup instrument. &
  Slow down. Use a hand lead or backup echo sounder. \\
Autopilot not responding &
  Tap \key{Stand By}. Take manual helm. &
  Investigate the cause before re-engaging. \\
Data stale everywhere &
  Check connection state. &
  Allow recovery or check Signal~K server. \\
Data missing for one value &
  Cross-check with another source. &
  Check the sensor and the Signal~K path. \\
\bottomrule
\end{longtable}

\section*{One-Tap Navigation}

\rowcolors{2}{LtGray}{white}
\begin{longtable}{L{5.8cm}L{7.9cm}}
\toprule
\rowcolor{Navy}
\color{white}\sffamily\bfseries Goal &
\color{white}\sffamily\bfseries Action \\
\midrule
Return to home screen              & Tap \iconlg{applications}~\key{Apps} --- always available \\
Switch to another open app         & Tap its icon in the open-apps strip \\
Change comfort preset              & \ui{Settings~> Comfort} \\
Check connection status            & Status area in Top Bar (right side) \\
Mark anchor position               & Tap \iconlg{anchor-iec}~\key{Anchor} \textrightarrow{} \iconlg{anchor-drop}~\key{Drop} \\
Cancel anchor watch                & Tap \iconlg{anchor-iec}~\key{Anchor} \textrightarrow{} \iconlg{anchor-raise}~\key{Raise} \\
Cancel a POB event                 & POB bar \textrightarrow{} select incident \textrightarrow{} \key{Cancel} \\
Export voyage track                & \ui{My~Data~> Tracks} \textrightarrow{} \key{Export} \\
\bottomrule
\end{longtable}

\chapter{Safety Notes}
\label{ch:safety}

\begin{warningbox}
FairWindSK is an aid to navigation.
It does not replace seamanship, lookout, paper or independently powered charts,
or any certified safety procedure.
The master and crew of the vessel remain solely responsible for the safe conduct of
the voyage.
\end{warningbox}

\begin{warningbox}
Always cross-check safety-critical information --- position, depth, heading,
collision risk --- with independent instruments before acting on it.
If FairWindSK contradicts what you see outside the boat or on trusted instruments,
trust your eyes and independent navigation references first.
\end{warningbox}

\begin{warningbox}
Missing depth is not zero depth.
A missing value means the measurement is unavailable.
Treat it as entirely unknown and navigate with full caution.
\end{warningbox}

\begin{warningbox}
After a reconnect or server restart, reconfirm all key navigation values before
depending on them.
Give GPS, depth, wind, and heading sensors time to re-acquire and settle.
\end{warningbox}

\begin{warningbox}
Autopilot commands must be confirmed against the vessel's actual behaviour.
Never assume a mode change took effect because the button responded.
Stand by to disengage and take manual helm at any moment.
\end{warningbox}

\begin{warningbox}
Activating an alarm in FairWindSK raises a Signal~K notification on the onboard
network.
It does not contact coast guard, MRCC, or any external authority.
Make a DSC and/or voice distress call on VHF Channel~16 in any real emergency.
\end{warningbox}

\begin{warningbox}
In an emergency, prioritise boat handling, crew safety, and direct observation over
any interaction with an electronic display.
A screen is a tool; your seamanship is your safety.
\end{warningbox}

\vfill
\begin{center}
  {\color{Teal}\rule{0.5\textwidth}{1pt}}\\[10pt]
  {\Large\sffamily\bfseries\color{Navy}Safe sailing.}\\[8pt]
  {\normalsize\sffamily\color{gray}
    The OpenFairWind Team \textbullet{}
    University of Naples 'Parthenope'}\\[4pt]
  {\small\sffamily\color{Teal}\url{https://github.com/OpenFairWind/FairWindSK}}
\end{center}
